\section{Sparse Readout and Causal Involvement of Goal Heads}
\label{sec:heads}
A compact set may suffice to expose subgoal information without containing every path used for native implementation. We therefore separate readout performance at different head budgets from interventions that restrict the model's access to task information.

\subsection{Small head sets support effective readout}
\label{sec:heads:sparsity}
\begin{figure}[t]
\centering
\includegraphics[width=.80\linewidth]{Figures/sparsity_eq.pdf}
\caption{Readout versus head budget at the assignment-operator pivot, before implementation (253 matched specifications; norm-relative additive injection, $\alpha=.1$). Lines show available head selections; stars mark the separately selected 122-head reference. Selections are not strictly nested, so connecting points guides comparison rather than isolating an effect of head count alone. Other capture positions appear in \Cref{app:sparsity_pivots}.}
\label{fig:sparsity}
\end{figure}
At the assignment pivot, twenty-five selected heads---approximately 0.98\% of the model's 2,560 query heads---achieve 93.7\% readout faithfulness and 90.5\% implementation transfer. The 122-head reference achieves 97.2\% and 95.7\%, respectively (\Cref{fig:sparsity}). Small selected sets thus recover much of the readout performance at this pre-implementation position. Larger sets do not yield a monotonic improvement. This does not identify a minimal circuit or imply that arbitrary sets of this size suffice. Results at the previous-line end, indentation, variable name, and first implementation token are in \Cref{app:sparsity_pivots}; comparisons with other injected sources appear in \Cref{app:readout_sources}.

\subsection{Preserving or restricting task access}
\label{sec:mechanism:ablation}
We evaluate native frontier implementation on the same 253 specifications at the assignment pivot, with interventions continuing during decoding. Each example has a clean baseline, a goal-head condition, and a random-head control. We test the core 122 heads and a broader 250-head set.

\paragraph{Conditional sufficiency: preserve selected access.}
For heads in layers 20 and above outside the spared set, we block attention to the span from just before the frontier description through the end of the prompt. Spared heads retain access. Other key positions remain available, layers below 20 are intact, and all MLP computation continues. Thus this intervention tests sufficiency of \emph{selected upper-layer access to the blocked span}, not sufficiency of these heads in isolation. Random controls spare a comparable number of heads from the specified depth pool; the pool and exact blocked span are detailed in \Cref{app:ablation}.

\paragraph{Necessity test: restrict selected heads to the sink.}
From the pivot onward, we instead restrict the selected heads to attend only to position zero, while other heads operate normally. This removes their normal content-dependent attention but is not literal output zeroing. We compare with restricting an equal-size random set drawn from the depth pool and disjoint from the goal heads.

\begin{table}[t]
\centering\small
\begin{tabular}{lcc}
\toprule
Condition & 122 heads & 250 heads\\
\midrule
Clean baseline & .925 & .925\\
Spare goal-head access & .806 & .913\\
Spare random-head access & .091 & .095\\
\midrule
Restrict goal heads to sink & .854 & .158\\
Restrict random heads to sink & .925 & .921\\
\bottomrule
\end{tabular}
\caption{Native frontier implementation under span restriction and sink-attention interventions ($n=253$). The two experiments manipulate different access patterns; ``spare'' does not remove the rest of the network.}
\label{tab:goal-ablation}
\end{table}

Preserving goal-head access retains substantially more implementation performance than preserving random-head access (\Cref{tab:goal-ablation}). Conversely, sink-restricting 250 goal heads sharply reduces success, whereas the random control remains close to baseline. Restricting the 122-head core has a smaller but nonzero effect. These results establish causal involvement of the broader selection under the tested intervention. They do not prove a unique bottleneck, identify the computation that first selects the goal, or demonstrate active rerouting after ablation.

The combination is informative: a small set supports effective readout, while native implementation remains robust to restricting the core and is much more vulnerable to restricting the broader set. Readout sufficiency and the extent of native dependence are different properties of the representation.
